% SOURCES: README.md; docs/1.0.0v/supervisor-report.md (§1, §13); draft/Features/features.md
% GUIDANCE: 150–250 words. State (1) the problem, (2) the approach, (3) what was
% built, (4) the headline results. Write this LAST, after ch. 11 is final.

% TODO: draft the abstract. Skeleton to fill:
Draw-to-Web is a desktop application that lets users author modern, semantic,
responsive, and accessible web pages by composing them on a canvas, with no code
required. % (1) problem: existing visual builders emit bloated, non-semantic markup.
% (2) approach: a single Document Model as the source of truth, a deterministic
%     generator, a tokens-driven styling system, and an accessibility hard gate.
% (3) built: Electron + React editor, nine-stage export pipeline, opt-in runtime
%     snippets, plus the v1.0.0 authoring surfaces (draw-to-create, match-layout,
%     and a headless MCP server).
% (4) results: <summarise the ch. 11 numbers — test suite, a11y gate, output/perf
%     budgets>. % TODO: fill from ch. 11 (do not invent the pending #95 numbers).

\vspace{0.5em}
\noindent\textbf{Keywords:} % TODO
visual web authoring, no-code, semantic HTML, accessibility, design tokens,
code generation, Electron, Model Context Protocol.
